\section{Related Work}
\label{sec:related}
We review five strands of work and state, for each, what is new here; no new optimizer is proposed.

\subsection{Wind Farm Layout Optimization with Metaheuristics}
\label{sec:related-wflop}
\emph{Discrete and continuous formulations.} Early studies placed turbines in the cells of a square grid and optimized the cell occupation with genetic algorithms (GAs)~\cite{Mosetti1994,Grady2005}; binary PSO with time-varying acceleration coefficients was applied to the same kind of grid~\cite{Pookpunt2013}. A grid makes the problem combinatorial but restricts the layouts that can be found. Kusiak and Song~\cite{Kusiak2010} formulated a continuous model with turbine coordinates in a circular farm, a discretized wind rose, a Jensen wake~\cite{Jensen1983,Katic1986} and explicit boundary and spacing constraints. This model, with two wind data sets~\cite{Kusiak2010,Bansal2017} and several farm radii, has become a common reference for metaheuristics and defines the benchmark used here (Section~\ref{sec:setup}). On it, or on close variants, ant colony optimization and particle filtering~\cite{Eroglu2012,Eroglu2013}, biogeography-based optimization~\cite{Bansal2017} and random search inside the feasible space~\cite{Feng2015} have been proposed, and GA, PSO, differential evolution (DE) and simulated annealing were compared on offshore layouts~\cite{Elkinton2008}; PSO has reduced wake losses in~\cite{Asaah2021}.

\emph{Newer metaheuristics and hybrids.} Later work used self-informed GAs~\cite{Ju2019}, an adaptive evolutionary algorithm combined with Monte Carlo tree search~\cite{Bai2022}, deterministic refinement techniques~\cite{Nagpal2021} and VNS for large offshore farms~\cite{Cazzaro2022}. Metaphor-based algorithms have also been applied; for example, the water cycle algorithm was compared with the SSA and the grey wolf optimizer on turbine placement~\cite{Rezk2019}, and the Laplacian SSA of two of the authors~\cite{Solanki2023} was combined with VNS in an earlier comparison that motivated this study. Reviews list many such methods, often evaluated on the same classical wind scenarios but with differing budgets, settings and constraint handling~\cite{Azlan2021}.

\emph{Common benchmarks.} The GECCO layout competitions provided a common evaluation interface, and the top entries of the 2015 edition were compared in~\cite{Wilson2018}. The IEA Wind Task~37 case studies fixed a simplified Gaussian wake, site boundaries and wind roses~\cite{Baker2019,IEA37repo}, and eight methods, gradient-based, gradient-free and hybrid, each run by researchers experienced with it, reached different layouts of similar energy yield on a common problem~\cite{Thomas2023}. Both efforts compare methods \emph{as submitted}: each team chose its own settings and effort, so the results show what experts achieve with their methods, not how a given choice of baseline setting or control changes a ranking.

\emph{Gradient-based and model-based alternatives.} Gradient-based optimization with exact gradients reduces the cost of layout optimization considerably compared with finite differences~\cite{Guirguis2016}; pseudo-gradients offer a cheap alternative to exact gradients~\cite{Quaeghebeur2021}, and reduced parameterizations of the layout cut the number of design variables~\cite{Stanley2019}. Gaussian wake models~\cite{Bastankhah2014,Tao2020} make the objective smoother than the top-hat Jensen wake, and surrogate models extend layout optimization to larger farms and higher-fidelity flow models~\cite{Yang2023,Wang2024,Li2025}.

\emph{Position of this study.} We use the Kusiak--Song benchmark in its continuous form (discretized only in the wind directions), and test beyond it on a Horns Rev~1 block and IEA37 Case Study~1 (Section~\ref{sec:beyond}). Unlike studies that propose a new method, we fix the methods and vary the configuration and the controls; unlike competitions, every method is run by the same code under one budget in objective evaluations, which for the gradient-based baseline (MS-SLSQP) includes its finite-difference gradients. We do not test gradient-based methods on the smooth models and with the exact gradients for which they are designed~\cite{Guirguis2016,Thomas2023}.

\subsection{Benchmarking Methodology for Metaheuristics}
\label{sec:related-bench}
Hooker~\cite{Hooker1995} distinguished competitive testing, which tells which algorithm wins, from controlled experiments, which tell why. Later guidelines turned this into practice: equal budgets counted in function evaluations, reported and justified settings of all methods, well-chosen reference algorithms and component analyses of new proposals~\cite{BartzBeielstein2020,LaTorre2021}. Platforms such as COCO~\cite{Hansen2021} standardize budgets, instances and performance measures for continuous black-box optimizers, but for artificial test functions; a constrained engineering problem such as the WFLOP must additionally decide how infeasible final solutions enter the comparison.

For the statistics, Dem\v{s}ar~\cite{Demsar2006} and Derrac et~al.~\cite{Derrac2011} recommended nonparametric tests over multiple problems: the Friedman test with post hoc comparisons of average ranks and pairwise Wilcoxon signed-rank tests, with a multiplicity correction such as Holm's~\cite{Holm1979}. Benavoli et~al.\ showed that post hoc tests based on mean ranks depend on the other methods in the pool~\cite{Benavoli2016}, which is why our primary tests use case means, and advocated Bayesian analyses with a region of practical equivalence~\cite{Benavoli2017}; Carrasco et~al.~\cite{Carrasco2020} reviewed these classical and newer trends for swarm and evolutionary algorithms and gave practical guidelines. Campelo and Takahashi~\cite{Campelo2019} derived the numbers of instances and runs needed for a given power and accuracy. Equivalence testing with two one-sided tests (TOST) asks whether a difference lies within a margin of practical irrelevance~\cite{Lakens2017}, and so avoids reading a nonsignificant difference as evidence that two algorithms perform equally.

\emph{Position of this study.} We apply this practice to the WFLOP and add three elements that we have not found reported for this problem: feasibility-aware ranking of penalty-based methods, equal-cost random-sampling controls for the first phase of hybrids, and equivalence analyses (TOST and a Bayesian signed-rank test) at three levels of generalization, over seeds, over cases and over farm clusters (Section~\ref{sec:setup}). We did not size the experiment by a power analysis as in~\cite{Campelo2019}; instead, for each equivalence statement, we report the smallest margin at which it holds, so that a reader can apply a different margin.

\subsection{Critiques of Metaphor-Based Algorithms}
\label{sec:related-metaphor}
S\"orensen~\cite{Sorensen2015} argued that many ``novel'' metaheuristics hide known search operators behind a new metaphor and are justified by comparisons with weak baselines; a group of editors and researchers later called for action on the proliferation of such methods~\cite{Aranha2022}. Camacho-Villal\'on et~al.~\cite{CamachoVillalon2023} showed for six popular metaphor-based algorithms that their search rules reuse known ideas under new names and are presented in a misleading way. For the SSA~\cite{Mirjalili2017}, the critical review of Castelli et~al.~\cite{Castelli2022} analyzed the update rules and the evidence for its performance; in the form used here, the leaders perturb the best solution uniformly with a shrinking radius and the followers only average (Section~\ref{sec:ssa}).

\emph{Position of this study.} These critiques rest mostly on the update equations or on artificial test functions. We add a component-level test on an applied constrained problem: whether a salp-swarm first phase does better than random sampling of equal cost in the feasible region. It does not (Section~\ref{sec:ablation}). This is also a self-critique, as LX-SSA~\cite{Solanki2023} was proposed by two of the authors; the test assesses LX-SSA as published, including its cost of $2N_p$ evaluations per iteration, and does not separate its Laplace step from that cost.

\subsection{PSO Parameterization and Stability Theory}
\label{sec:related-pso}
The original PSO~\cite{Kennedy1995} used acceleration coefficients $c_1=c_2=2$ and velocity clamping; the inertia weight~\cite{Shi1998} and the constriction factor of Clerc and Kennedy~\cite{Clerc2002} were introduced to control the swarm's divergence, and the constriction setting is equivalent to an inertia-weight PSO with $w\approx0.73$ and $c_1=c_2\approx1.50$. Analyses of the particle dynamics~\cite{Trelea2003,VandenBergh2006} derived the parameter regions in which a deterministic simplification of a particle converges and related the parameters to the exploration behavior. Poli~\cite{Poli2009} derived the mean and variance of the sampling distribution under stagnation and thereby the order-1 and order-2 stability regions; Cleghorn and Engelbrecht~\cite{Cleghorn2018} showed that, for $c_1=c_2$, the order-2 region also holds under a weaker, non-stagnant assumption; and Bonyadi and Michalewicz~\cite{Bonyadi2017} reviewed these results together with the many PSO variants and their parameter choices.

\emph{Position of this study.} These results are well known in PSO theory but rarely used when PSO serves as a baseline in applied comparisons. We use the order-2 condition as a check of a baseline setting: the setting $w=0.7$, $c_1=c_2=2$ is order-1 but not order-2 stable (Proposition~\ref{prop:pso-stability}), a sweep of $c_1=c_2$ shows that feasibility on the benchmark begins to degrade near the boundary, and the setting changes the ranking of PSO against the salp-swarm hybrids (Section~\ref{sec:psosetting}). The link between a stability region of the unconstrained dynamics and the feasibility of the returned solutions in a constrained problem is, to our knowledge, not reported for the WFLOP.
% CHECK-THEORY [T01]/[T04]: order-2 violated (4 > 3.50); CHECK-CSWEEP [S02]-[S05]: degradation begins near c*; CHECK-FINAL [C27]: old setting reverses PSO vs salp-swarm hybrids

\subsection{Hybrid Swarm--Local-Search Designs}
\label{sec:related-hybrid}
Combining a population-based global search with a local search is the idea of memetic algorithms; Krasnogor and Smith~\cite{Krasnogor2005} gave a taxonomy and discussed design issues such as the choice of the local search and its integration into the evolutionary cycle, and Neri and Cotta~\cite{Neri2012} reviewed memetic computing for discrete, continuous and constrained problems. VNS~\cite{Mladenovic1997,Hansen2001} changes neighborhoods systematically around an incumbent and has a continuous form~\cite{Mladenovic2008}; in the WFLOP, it has been used for large offshore farms~\cite{Cazzaro2022}, and deterministic refinement techniques for layouts have been compared in~\cite{Nagpal2021}. The simplest of these designs is the sequential (relay) hybrid, which runs the global method for part of the budget and passes its best solution to the local method.

\emph{Position of this study.} PSO-VNS is such a sequential hybrid, built from a standard PSO and the basic VNS without new operators. What is new is its evaluation: the first phase is compared with equal-cost random sampling in the farm and in its bounding square, and with the local search alone, so that the contribution of each phase is measured rather than assumed; and the second phase is instrumented, which shows that at 6,030 evaluations the VNS acts essentially as one compass-search descent, so the comparison measures the value of a start for a local search (Section~\ref{sec:ablation}).
% CHECK-FINAL [C53]-[C56]: component analysis verdicts (Section ablation); CHECK-DIAG [D08]-[D10], [D20]: VNS phase essentially one compass-search descent
